% SGA 3, Expose IV, Proposition 4.5.2 through Corollary 4.5.3.
% French authority: Exp4-13oct24.pdf, local pages 44--45
% (combined-reader pages 232--233; printed pages 220--221).
% The external English Markdown was used only as a comparison.

\medskip
\noindent\SGARefTarget{sga3:iv:proposition:4:5:2}\textbf{命题 4.5.2.}
\textit{设 \(S\) 是 \(\mathscr C\) 上的一个可表层,
\(F\to S\) 是 \(\widehat{\mathscr C}\) 中的态射.
\(\mathscr C/S\) 上的预层
\[
  S'\longmapsto \operatorname{Hom}_{S}(S',F)
\]
是分离的 (分别地, 是层), 当且仅当 \(F\) 是 \(\mathscr C\)
上的分离预层 (分别地, 层).}

对每个函子 \(P\), 有 (I, \SGARefLink{sga3:i:proposition:1:4:1}{1.4.1})
\[
  \operatorname{Hom}(P,F)
  =\coprod_{f\in\operatorname{Hom}(P,S)}
    \operatorname{Hom}_{f}(P,F).
\]
\footnote{编者注: 下文扩充了原文中隐含的论证.}
设 \(S'\in\operatorname{Ob}(\mathscr C)\),
\(\eta:C'\to S'\) 是一个覆盖筛.
由于 \(S\) 是层, 映射 \(f\mapsto f\circ\eta\) 是双射
\[
  \operatorname{Hom}(S',S)\xrightarrow{\sim}
  \operatorname{Hom}(C',S).
\]
所以映射
\[
  \operatorname{Hom}_{\widehat{\mathscr C}}(S',F)
  \longrightarrow
  \operatorname{Hom}_{\widehat{\mathscr C}}(C',F)
\]
分解为映射
\[
  \operatorname{Hom}_{f}(S',F)
  \longrightarrow
  \operatorname{Hom}_{f\circ\eta}(C',F).
\]
的 ``不交并''. 由此即得\SGARefLink{sga3:iv:proposition:4:5:2}{此命题}.

\medskip
\noindent\SGARefTarget{sga3:iv:corollary:4:5:3}\textbf{推论 4.5.3.}
\textit{若 \(\mathscr C\) 上的拓扑比典范拓扑更粗,
则它在 \(\mathscr C/S\) 上诱导的拓扑,
比 \(\mathscr C/S\) 的典范拓扑更粗.}

% Translation coverage: entire assigned source component, without omissions.
% Statistics: source lines=49; target lines=52; inline math=15; display math=5; unresolved translation TODOs=0.
